\section{Proof of Proposition~\ref{prop:aging-expansion}}\label{app:aging}

We bound a gamma ratio, bracket $C_N$ and $c_N$, prove a local limit theorem
at the center, and then expand the crossing. We use the setting of
Sections~\ref{sec:notation} and~\ref{sec:aging} with 2 states and a fair first step. So $0<c<2$, $X_1$ is a fair sign, and the switch indicators
$\chi_k=\mathbf 1\{X_k\ne X_{k-1}\}$, $k\ge2$, are independent
Bernoulli$(c/k)$ variables, independent of $X_1$. Put
\[
\rho(a,b)=\prod_{k=a+1}^{b}\Bigl(1-\frac ck\Bigr)
=\frac{\Gamma(b+1-c)\,\Gamma(a+1)}{\Gamma(b+1)\,\Gamma(a+1-c)},\qquad
\lambda_{t,N}=\frac ct\,\rho(t,N)
\]
for integers $1\le a\le b$ and $2\le t\le N$. So $\lambda_{t,N}$ is the
probability of a switch at time $t$ and none in $(t,N]$, and
$E_N=\frac12\rho(1,N)$. Let $T^+_s$ be the time of the $s$-th step equal to
$+1$. Let $B$ be a Beta$(c,c)$ variable, $d_c=4^{1-c}/\mathrm B(c,c)$ its
density at $\frac12$, and $\Theta(c)=\log\bigl(2\Gamma(2-c)\,d_c\bigr)$. For
even $N$ we write $h=N/2$ and $C_N=\Prob(A_N=h)$. Changing the sign of every
step preserves the law of the walk and sends $A_N$ to $N-A_N$. We call this
flip symmetry.

\subsection{A gamma ratio}

Gautschi's inequality~\cite[5.6.4]{dlmf} says that
$x^{1-s}<\Gamma(x+1)/\Gamma(x+s)<(x+1)^{1-s}$ for $x>0$ and $0<s<1$. We
claim that for $0<c<2$ and real $y\ge2$,
\begin{equation}\label{eq:D-gamma}
\frac{\Gamma(y+1-c)}{\Gamma(y+1)}=y^{-c}(1+\vartheta)\qquad\text{with}\qquad
-\frac1y\le\vartheta\le\frac2y .
\end{equation}
If $c=1$ the left side is $1/y$ and $\vartheta=0$. If $c<1$ we take $x=y$ and
$s=1-c$. Then the left side lies between $(y+1)^{-c}$ and $y^{-c}$, and
$(1+1/y)^{-c}\ge1-c/y\ge1-1/y$. If $c>1$ we take $x=y$ and $s=2-c$ and use
$\Gamma(y+1-c)=\Gamma(y+2-c)/(y+1-c)$. Then the left side lies between
$(y+1)^{1-c}/(y+1-c)$ and $y^{1-c}/(y+1-c)$. Here
$y/(y+1-c)=1+(c-1)/(y+1-c)\le1+1/(y-1)\le1+2/y$, and
$(1+1/y)^{1-c}\,y/(y+1-c)\ge(1+1/y)^{-1}\ge1-1/y$.

\subsection{Averaging at the center and a bracket}

\begin{lemma}[averaging at the center]\label{lem:D-average}
Let $N\ge2$ be even and $T=T^+_h$. Then
\begin{equation}\label{eq:D-average}
C_N=2\sum_{r=h}^{N-1}\Prob(T=r)\,\lambda_{r+1,N}\qquad\text{and}\qquad
\sum_{r=h}^{N-1}\Prob(T=r)=\frac12 .
\end{equation}
If $0<c\le1$, then $c/N\le C_N\le\lambda_{h+1,N}$.
\end{lemma}

\begin{proof}
On the event $\{A_N=h,\,X_N=-1\}$ the last step equal to $+1$ is the $h$-th
one, at a time $T=r\le N-1$. It is followed by a switch at $r+1$ and no switch
in $(r+1,N]$. Conversely, these conditions give $A_N=h$ and $X_N=-1$. The
event $\{T=r\}$ depends only on $X_1,\chi_2,\dots,\chi_r$. So
$\Prob(A_N=h,\,X_N=-1)=\sum_{r=h}^{N-1}\Prob(T=r)\,\lambda_{r+1,N}$, and flip
symmetry gives the same value for $X_N=+1$. Next, $T\le N-1$ exactly when
$A_{N-1}\ge h$. This has probability $\frac12$, because $N-1$ is odd and
$A_{N-1}$ is symmetric about $(N-1)/2$. Finally
$\lambda_{t+1,N}/\lambda_{t,N}=t/(t+1-c)$, which is at most $1$ for $c\le1$.
So $\lambda_{N,N}=c/N\le\lambda_{t,N}\le\lambda_{h+1,N}$ for
$h+1\le t\le N$, and by~\eqref{eq:D-average} $C_N$ is a weighted average of
these values.
\end{proof}

For an integer $M\ge2$ put $\eta_M(c)=M\rho(1,M)/(2c)$. Since $\rho(1,M)$
decreases in $c$ and equals $1/M$ at $c=1$, the function $\eta_M$ decreases
strictly on $(0,1]$, from $\infty$ at $0$ to $\eta_M(1)=\frac12$. Let
$c^{(M)}\in(0,1)$ be the root of $\eta_M=1$.

\begin{proposition}[bracket]\label{prop:D-bracket}
The numbers $c^{(M)}$ increase with $M$, and
$1-c^{(M)}=\frac{\log2}{\log M+1+\gamma}+O\bigl((\log M)^{-3}\bigr)$. For
every even $N\ge2$ we have $c^{(h+1)}\le c_N\le c^{(N)}$. So
$(1-c_N)\log N\to\log2$, and $c_N\ge c^{(3)}=(9-\sqrt{57})/2>0.725$ for
$N\ge4$.
\end{proposition}

\begin{proof}
The equation $\eta_M(c)=1$ says that $\frac12\rho(1,M)=c/M$, namely that the
end for $M$ steps equals $c/M$. We compare $E_N$ with the 2 bounds for $C_N$
in Lemma~\ref{lem:D-average}, and each comparison is an inequality for some
$\eta_M$.

\emph{Monotonicity.} For $c<1$ we have $\eta_{M+1}/\eta_M=(M+1-c)/M>1$. So
$\eta_{M+1}(c^{(M)})>1$, and since $\eta_{M+1}$ decreases, $c^{(M+1)}>c^{(M)}$.

\emph{The bracket.} By Proposition~\ref{prop:aging-fair}(b), $C_N<E_N$ for
$c<c_N$ and $C_N>E_N$ for $c>c_N$, and $c_N<1$. So it is enough to compare
$C_N$ with $E_N$ for $0<c<1$. Let $c^{(N)}<c<1$. Then $\eta_N(c)<1$, that is
$E_N=\frac12\rho(1,N)<c/N$, and $c/N\le C_N$ by Lemma~\ref{lem:D-average}. So
$C_N>E_N$, which means $c>c_N$. Hence $c_N\le c^{(N)}$. Now let
$c<c^{(h+1)}$. Then $\eta_{h+1}(c)>1$, that is
$\frac12\rho(1,h+1)>c/(h+1)$. With $\rho(1,N)=\rho(1,h+1)\rho(h+1,N)$ this
gives $E_N>\frac{c}{h+1}\rho(h+1,N)=\lambda_{h+1,N}\ge C_N$. So $c<c_N$.
Hence $c_N\ge c^{(h+1)}$.

\emph{The expansion.} The number $c^{(3)}$ solves
$3\rho(1,3)=(2-c)(3-c)/2=2c$, that is $c^2-9c+6=0$. For $N\ge4$ we have
$h+1\ge3$, so $c_N\ge c^{(h+1)}\ge c^{(3)}$. By~\eqref{eq:D-gamma} with $y=M$
we have $M\rho(1,M)=M^{1-c}(1+O(1/M))/\Gamma(2-c)$. So the equation
$\eta_M(c)=1$ reads $(1-c)\log M=\Theta_0(c)+O(1/M)$ with
$\Theta_0(c)=\log\bigl(2c\,\Gamma(2-c)\bigr)$, uniformly for $0<c\le1$. For
$M\ge3$ we have $c^{(M)}\in[c^{(3)},1)$, where $\Theta_0$ is bounded. So
$\varepsilon=1-c^{(M)}$ is $O(1/\log M)$, and we may expand at $c=1$. We have
$\Theta_0(1)=\log2$ and $\Theta_0'(1)=1-\psi(1)=1+\gamma$, where
$\psi=\Gamma'/\Gamma$ and $\psi(1)=-\gamma$~\cite[5.4.12]{dlmf}. Taylor's
theorem gives
$\varepsilon(\log M+1+\gamma)=\log2+O(\varepsilon^2)+O(1/M)$, which is the
expansion. The limit
of $(1-c_N)\log N$ follows since $\log(h+1)\sim\log N$.
\end{proof}

\subsection{A local limit at the center}

\begin{proposition}[local limit at the center]\label{prop:D-llt}
For each compact $J\subset(0,2)$ there is $C_J$ such that
$\bigl\lvert NC_N-d_c\bigr\rvert\le C_J\sqrt{\log N/N}$ for $c\in J$
and even $N\ge4$.
\end{proposition}

\begin{proof}
We start from the averaging identity of Lemma~\ref{lem:D-average} and replace
the law of $T=T^+_h$ by its continuum version. Since $T\le r$ exactly when
$A_r\ge h$, we need the law of $A_r$. The chain of
Proposition~\ref{prop:aging-urn} with $d=2$ and $\theta=c$ switches with
nearly the same probabilities as the walk, and its counts are binomial with a
random parameter $B$. So $A_r$ is close in law to $rB$, and the time $T$ is
close in law to $h/B$. The weights $\lambda_{r+1,N}$ are explicit, so the
average becomes an explicit integral against the Beta$(c,c)$ density.

The proof has 5 steps. Step~1 shows that the counts of the chain are
binomial with the random parameter $B$. Step~2 couples the walk with the chain
and bounds the difference of the distribution functions of $T$ and $h/B$. Step~3 uses
summation by parts to replace the law of $T$ by that of $h/B$ in the average.
Step~4 writes the new sum as an expectation over $B$, and Step~5 evaluates it.

Let $F(r)=\Prob(T\le r)=\Prob(A_r\ge h)$, $\Phi(r)=\Prob(rB\ge h)$ and
$\Lambda_r=N\lambda_{r+1,N}$. By Lemma~\ref{lem:D-average},
\[
NC_N=2\sum_{r=h}^{N-1}\bigl(F(r)-F(r-1)\bigr)\Lambda_r .
\]
In this proof $C$ depends only on $J$ and may change from line to line. Since
$NC_N\le N$ and $d_c$ is bounded on $J$, it is enough to prove the claim for
large $N$. We put $K=\lceil\sqrt N\rceil$, $u=\sqrt{N\log N}$ and $w=K+u$, and
we take $N$ so large that $w\le h/4$. Note that $K\le h$.

\emph{Step 1 (a chain with binomial counts).} Let $X'$ be the chain of
Proposition~\ref{prop:aging-urn} with $d=2$ and $\theta=c$, written with the
states $\pm1$, and let $A'_r$ be its number of steps $+1$ among the first $r$.
Its first step is a fair sign, and at step $k\ge2$ it switches with
probability $q_k=c/(k-1+2c)$, independently of the past. By
Proposition~\ref{prop:aging-urn} and $x^{(m)}=\Gamma(x+m)/\Gamma(x)$,
\[
\Prob(A'_r=s)=\binom rs\frac{c^{(s)}c^{(r-s)}}{(2c)^{(r)}}
=\binom rs\frac{\mathrm B(c+s,c+r-s)}{\mathrm B(c,c)}
=\E\Bigl[\binom rsB^s(1-B)^{r-s}\Bigr].
\]
So $A'_r$ has the law of a binomial variable with $r$ trials and the random
success probability $B$. By Hoeffding's inequality~\cite[Theorem~1]{hoeffding1963},
a binomial variable with $r$ trials and success probability $x$ is at least
$rx+u$ with probability at most $e^{-2u^2/r}$, and the same holds for being at
most $rx-u$. Hence for every real $y$
\[
\Prob(rB\ge y+u)-e^{-2u^2/r}\le\Prob(A'_r\ge y)\le\Prob(rB\ge y-u)+e^{-2u^2/r}.
\]

\emph{Step 2 (coupling the walk with $X'$).} The switch probabilities differ
by $c/k-q_k=c(2c-1)/(k(k-1+2c))$, which is at most $6/(k(k-1))$ in absolute
value since $c<2$. Let $K\le r\le N$. For $K<k\le r$ we couple the switch
indicator $\chi_k$ of the walk with the switch indicator $\chi'_k$ of $X'$ so
that $\Prob(\chi_k\ne\chi'_k)=\lvert c/k-q_k\rvert$, using independent extra
randomness. We put $X_K=X'_K$ and $X_k=X_{k-1}(-1)^{\chi_k}$ for $K<k\le r$,
and we draw $X_1,\dots,X_{K-1}$ from the law of the walk given $X_K$, again
with independent randomness. In both processes the sign at step $K$ is fair,
and the later switches are independent of the first $K$ steps. So
$(X_1,\dots,X_r)$ has the law of the walk. Outside an event of probability at
most $\sum_{k>K}6/(k(k-1))=6/K$ the two processes agree at the steps
$K,\dots,r$, and then $\lvert A_r-A'_r\rvert<K$. With Step 1, $w=K+u$ and
$e^{-2u^2/r}\le e^{-2\log N}$, this gives
\[
\lvert F(r)-\Phi(r)\rvert\le\Prob(h-w\le rB<h+w)+\frac1{N^2}+\frac6K
\qquad(K\le r\le N).
\]

\emph{Step 3 (replacing $F$ by $\Phi$).} Put $\Delta(r)=F(r)-\Phi(r)$. Then
$\Delta(h-1)=0$, because $A_{h-1}<h$ and $B<1$. Summation by parts gives
\[
\sum_{r=h}^{N-1}\bigl(\Delta(r)-\Delta(r-1)\bigr)\Lambda_r
=\Delta(N-1)\Lambda_{N-1}-\sum_{r=h}^{N-2}\Delta(r)\bigl(\Lambda_{r+1}-\Lambda_r\bigr).
\]
The ratio $\lambda_{t+1,N}/\lambda_{t,N}=t/(t+1-c)$ gives
$\Lambda_{r+1}-\Lambda_r=\Lambda_r(c-1)/(r+2-c)$. Every factor of
$\rho(r+1,N)$ lies in $(0,1)$, so $\Lambda_r=\frac{Nc}{r+1}\rho(r+1,N)\le2c$
for $r\ge h$. Hence $\Lambda_{N-1}=c$ and
$\lvert\Lambda_{r+1}-\Lambda_r\rvert\le2c/h\le8/N$, and the right side is at
most $c\lvert\Delta(N-1)\rvert+\frac8N\sum_{r=h}^{N-2}\lvert\Delta(r)\rvert$
in absolute value. For each value $b$ of $B$ the number $rb$ lies in
$[h-w,h+w)$ for at most $2w/b+1$ integers $r$, and for no $r\le N$ unless
$b\ge(h-w)/N\ge\frac38$. So the probabilities in Step 2 add up to at most
$6w+1$ over $r$, and
$\sum_{r=h}^{N-2}\lvert\Delta(r)\rvert\le6w+2+3\sqrt N$. For $r=N-1$ the
event in Step 2 puts $B$ in an interval of length $2w/(N-1)$ inside
$[\frac38,\frac34]$, where the density of $B$ is at most $C$ for $c\in J$. So
$\lvert\Delta(N-1)\rvert\le Cw/N+N^{-2}+6/K$. Together, replacing $F$ by
$\Phi$ changes $NC_N$ by at most $C(w/N+1/K)\le C\sqrt{\log N/N}$.

\emph{Step 4 (the new sum is an expectation).} After the replacement the sum
is
\[
 2\sum_{r=h}^{N-1}\bigl(\Phi(r)-\Phi(r-1)\bigr)\Lambda_r .
\]
The term with $r=h$
vanishes, since $\Phi(h)=\Phi(h-1)=0$. For $h<r\le N-1$ we have
$\Phi(r)-\Phi(r-1)=\Prob(h/r\le B<h/(r-1))$, and on this event
$\lceil h/B\rceil=r$. These events are disjoint, and their union is
$\{h/(N-1)\le B<1\}$. So the new sum is
$2\,\E\bigl[\Lambda_{\lceil h/B\rceil};\,B\ge h/(N-1)\bigr]$.

\emph{Step 5 (evaluating the expectation).} By~\eqref{eq:D-gamma} with $y=N$
and $y=r+1\ge3$,
$\Lambda_r=c\bigl((r+1)/N\bigr)^{c-1}(1+O(1/N))$ for $h\le r\le N-1$. On the
event above with $r=\lceil h/B\rceil$, the number $(r+1)/N$ lies in
$[\frac12,2]$ and within $2/N$ of $h/(BN)=1/(2B)$, which also lies in
$[\frac12,2]$. The function $x^{c-1}$ is $4$-Lipschitz on $[\frac12,2]$. So
$\Lambda_{\lceil h/B\rceil}=c(2B)^{1-c}+O(1/N)$. Also
$\Prob\bigl(\frac12\le B<h/(N-1)\bigr)=O_J(1/N)$, since the density of $B$ is
bounded near $\frac12$ uniformly on $J$. So we may replace the event
$\{B\ge h/(N-1)\}$ by $\{B\ge\frac12\}$. In the remaining integral the factor
$b^{1-c}$ cancels the factor $b^{c-1}$ of the density. Hence, by Step 3,
$NC_N$ is within $C\sqrt{\log N/N}$ of
\[
2c\,\E\bigl[(2B)^{1-c};\,B\ge\tfrac12\bigr]+O_J(1/N)
=\frac{2c\,2^{1-c}}{\mathrm B(c,c)}\int_{1/2}^1(1-b)^{c-1}\,db+O_J(1/N)=d_c+O_J(1/N).\qedhere
\]
\end{proof}

\subsection{The expansion of the crossing}

\begin{proof}[Proof of Proposition~\ref{prop:aging-expansion}]
We write $\log C_N-\log E_N$ as an explicit function of $c$ with a small
error, and we solve for its zero near $c=1$. Let $J=[c^{(3)},1]$. Then $d_c$ is
bounded below on $J$ by a positive constant. By~\eqref{eq:D-gamma} with
$y=N$,
$\log E_N=\log\bigl(\tfrac12\rho(1,N)\bigr)=-\log2-\log\Gamma(2-c)-c\log N+O(1/N)$.
By Proposition~\ref{prop:D-llt},
$\log C_N=\log d_c-\log N+O(\sqrt{\log N/N})$ for large $N$. Both hold uniformly for
$c\in J$, so
\begin{equation}\label{eq:D-logratio}
\log C_N-\log E_N=(c-1)\log N+\Theta(c)+O(\sqrt{\log N/N})\qquad(c\in J).
\end{equation}
Here $\Theta(1)=\log2$ because $d_1=1$. Next,
$\Theta'(c)=-\psi(2-c)-\log4-2\psi(c)+2\psi(2c)$, and with $\psi(1)=-\gamma$
and $\psi(2)=1-\gamma$~\cite[5.4.12, 5.4.14]{dlmf} this gives
$\Theta'(1)=\gamma+2-2\log2$. Put
$L=\log N+\Theta'(1)$. By Proposition~\ref{prop:D-bracket}, $c_N\in J$ for
$N\ge4$ and $\varepsilon=1-c_N$ is $O(1/\log N)$. The left side
of~\eqref{eq:D-logratio} vanishes at $c=c_N$, because $C_N=E_N$ there. As
$\Theta$ is smooth on $J$, Taylor's theorem gives
$\Theta(c_N)=\log2-\varepsilon\,\Theta'(1)+O(\varepsilon^2)$.
So~\eqref{eq:D-logratio} at $c=c_N$ reads
$0=-\varepsilon\log N+\log2-\varepsilon\,\Theta'(1)+O(\varepsilon^2)+O(\sqrt{\log N/N})$,
that is
$\varepsilon L=\log2+O(\varepsilon^2)+O(\sqrt{\log N/N})=\log2+O(L^{-2})$. So
$\varepsilon=\log2/L+O(L^{-3})$, which is the expansion.
\end{proof}

With one more term in Taylor's theorem and
$\varepsilon^2=(\log2)^2/L^2+O(L^{-4})$, the same argument adds the term
$-\Theta''(1)(\log2)^2/(2L^3)$ to the expansion of $c_N$, with error
$O(L^{-4})$. Here $\Theta''(1)=\pi^2/2-4$, since
$\Theta''(c)=\psi'(2-c)-2\psi'(c)+4\psi'(2c)$, $\psi'(1)=\zeta(2)$ and
$\psi'(2)=\zeta(2)-1$~\cite[5.4.12, 5.15.5]{dlmf}. At $N=200$ the bracket of
Proposition~\ref{prop:D-bracket} is $[0.88749,0.89878]$,
Proposition~\ref{prop:aging-expansion} gives $0.89319$, and with the $L^{-3}$
term we get $0.89236$. The high-precision value is $c_{200}=0.89253$. The
constants $C_J$ and the constants in the $O$-terms of this appendix are not
computed. Numerically $\lvert NC_N-d_c\rvert$ decays like $1/N$. At $c=0.75$
it drops by a factor of about $4$ each time $N$ is multiplied by $4$, from
$N=400$ to $N=6400$. So the rate in Proposition~\ref{prop:D-llt} is probably
not sharp.
